\subsection{Binary trees and minimal orbits} \label{subsec:binarytrees}

In this section we reprove some important facts about binary trees.

\begin{definition}
A \emph{binary tree} is a rooted tree where each vertex has a left child, a right child, both children, or no children.
\end{definition}

Let $\mathcal{T}_n$ denote the set of binary trees on $n$ vertices. One of the interpretations of Catalan numbers is that $C_n=\lvert \mathcal{T}_n\rvert$. There is a symmetry group $G_n$ acting on $\mathcal{T}_n$, generated by reflections about a vertex which switch the left and right subtrees. An orbit of trees refers to an orbit under the action of $G_n$, which groups together binary trees which are isomorphic when the distinction between left and right is disregarded.

The following lemma was shown by Deutsch and Sagan \cite{deutsch2006congruences}.
\begin{lemma}
\label{lem:orbit sizes}
Let $\mathcal{O}$ be an orbit of the symmetry group $G_n$ on $\mathcal{T}_n$. Then $\lvert \mathcal{O}\rvert$ is of the form $2^t$ for a nonnegative integer $t$. Furthermore, $t\geq s=s_2(n+1)-1$, with equality for $(2s-1)!!$ orbits.
\end{lemma}

We will use \emph{minimal orbit} to refer to an orbit of the minimum possible size, which is $2^s$. The set of all orbits of $\mathcal{T}_n$ under $G_n$ is denoted $\mathcal{U}_n$, and the set of minimal orbits is denoted $\mathcal{U}^{\text{min}}_n$.

\begin{definition}
The \emph{complete binary tree of depth $k$} is the binary tree of $2^k-1$ vertices where each vertex in layer $0, 1, \dotsc, k-2$ has two children and each vertex in layer $k-1$ has no children.
\end{definition}
Note that complete binary trees are exactly the trees whose orbit consist of just one tree.

Now we have the following key result which describes the structure of minimal orbits of trees on $n$ vertices. (Recall that these are orbits of size $2^{s_2(n+1)-1}$.) This was done more generally for $q$-ary trees by Konvalinka \cite{konvalinka2007divisibility}.

\begin{theorem}[Structure of minimal orbits]
\label{thm:orbit structure}
Let $s=s(n+1)-1$ and $n+1=2^{k_1}+\dotsb+2^{k_{s+1}}$ be the binary expansion of $n+1$. Then all trees in an orbit of $\mathcal{U}_n^{\text{min}}$ may be constructed in the following manner: construct an arbitrary binary tree with $s$ vertices, and then attach completely symmetric trees of depth $k_1, \dotsc, k_{s+1}$ to the $s+1$ endpoints.
\end{theorem}

See Figure \ref{fig:minorbit} for an example. In constructing these minimal orbits, the $s$ vertices of the ``arbitrary binary tree'' will be marked black while the other vertices will be marked white.

\begin{figure}[h!]
\centering
    \begin{tikzpicture}[
    level distance=1cm,
    level 1/.style={sibling distance=4cm},
    level 2/.style={sibling distance=3cm},
    level 3/.style={sibling distance=1.5cm},
    level 4/.style={sibling distance=0.75cm},
    level 5/.style={sibling distance=0.375cm},
    level 6/.style={sibling distance=0.1875cm}]
      \node[black] {}
        child {node[black] {}
          child {node[white] {}
            child {node[white] {}
            child {node[white] {}
            child {node[white] {}
            child {node[white] {}}
            child {node[white] {}}}
            child {node[white] {}
            child {node[white] {}}
            child {node[white] {}}}
            }
            child {node[white] {}
            child {node[white] {}
            child {node[white] {}}
            child {node[white] {}}}
            child {node[white] {}
            child {node[white] {}}
            child {node[white] {}}}
            }
            }
            child {node[white] {}
            child {node[white] {}
            child {node[white] {}
            child {node[white] {}}
            child {node[white] {}}}
            child {node[white] {}
            child {node[white] {}}
            child {node[white] {}}}
            }
            child {node[white] {}
            child {node[white] {}
            child {node[white] {}}
            child {node[white] {}}}
            child {node[white] {}
            child {node[white] {}}
            child {node[white] {}}}
            }
            }
            }
          child {node[white] {}
          child {node[white] {}}
          child {node[white] {}}}
        }
        child {node[black] {}
        child {node[black, xshift=1.5cm] {}
        child {node[white] {}}
        child {node[white] {}
            child {node[white] {}
            child {node[white] {}
            child {node[white] {}}
            child {node[white] {}}}
            child {node[white] {}
            child {node[white] {}}
            child {node[white] {}}}
            }
            child {node[white] {}
            child {node[white] {}
            child {node[white] {}}
            child {node[white] {}}}
            child {node[white] {}
            child {node[white] {}}
            child {node[white] {}}}
            }
            }}
        };
    \end{tikzpicture}
\caption{An example for $n=54, n+1=2^5+2^4+2^2+2^1+2^0$.}
\label{fig:minorbit}
\end{figure}
Furthermore, this mapping can be reversed in the following way: for every vertex, color it white if the left and right subtrees at that vertex are isomorphic, and otherwise color it black. The black vertices are the original $s$ vertices and all white vertices are part of complete binary trees.

The symmetries of such trees are generated precisely by reflections about the original $s$ vertices (the ones marked in black).

\begin{remark}
In Theorem \ref{thm:orbit structure}, the term $2^0$ may appear in the binary expansion of $n+1$. This corresponds to an empty tree. In the example of Figure \ref{fig:minorbit}, all the black vertices have two children except for the third one from the left, which has only a right child. This is because the empty tree was assigned to its left child.
\end{remark}